\chapter{Einleitung}
\label{chap:introduction}

Mit der Veröffentlichung von ChatGPT 3.5 im November 2022 \autocite{openai2022} wurde der breiten Öffentlichkeit erstmals der Zugang zu einem niederschwelligen KI-Tool zur leichten Informationsbeschaffung ermöglicht, wo dies zuvor über Suchmaschinen erarbeitet werden musste. Schätzungen zufolge erreichte ChatGPT 3.5 innerhalb von zwei Monaten 100 Millionen Nutzer \autocite{hu2023}. Die steigende Nutzung und Nachfrage im beruflichen und privaten Alltag führt unweigerlich zu neuen Akteuren im Markt, darunter Google, Meta und Anthropic \autocite{le2025llmspredictprerequisiteskills}. Das Resultat ist eine enorme Menge an KI-generierten Inhalten in aller Form, also Text, Video, Bild und Audio. Die Art und Weise, wie digitale Inhalte konsumiert und reproduziert werden, hat sich folglich radikal verändert. Insbesondere Betreiber und Nutzer von sozialen Medien werden mit Inhalten konfrontiert, die durch Bots oder Sprachmodelle erzeugt werden. Daraus resultiert unweigerlich auch eine zunehmende und zunehmend einfachere Verbreitung von Mis- und Desinformation auf diesen Plattformen. Dies wirft gesellschaftliche und ethische Fragen auf, da Sprachmodelle gezielt zur Meinungsbeeinflussung genutzt werden können, indem authentisch und menschlich klingende Inhalte produziert und bewusst verbreitet werden \autocite{coeckelbergh2025}. \\
Wie so eine gezielte Meinungsbeeinflussung aussehen kann, hat sich im April 2025 gezeigt, als bekannt wurde, dass Forschende der Universität Zürich ein von Reddit bzw. dem Subreddit „r/changemyview“ nicht autorisiertes Experiment durchführten \autocite{redditmodcmv2025}. Im Rahmen des Experiments sollte herausgefunden werden, inwieweit sich Reddit-Nutzer von KI-Chatbots in ihrer Meinung beeinflussen lassen \autocite{osfunizurich}. Auf dem Subreddit „r/changemyview“ können Nutzer kontroverse Meinungen in Posts äußern und laden andere Nutzer damit zu einer argumentativen Diskussion ein, die die Meinung entweder verstärken oder den Nutzer zu einer Meinungsänderung bewegen können. Die Gegenstände der Diskussionen sind dabei nicht nur politisch-gesellschaftlicher Natur, sondern befassen sich auch mit alltäglichen Lebensfragen. Jeder Subreddit hat üblicherweise eine Reihe von Regeln, die dem Subreddit eigen sind und deren Einhaltung von einem Moderationsteam überwacht wird. Der Subreddit „r/changemyview“ verbietet in seinem Regelwerk explizit die Nutzung und Veröffentlichung KI-generierter Inhalte \autocite{cmvwiki}. Im Experiment der Universität Zürich wurden unter Zuhilfenahme von Sprachmodellen Kommentare erzeugt, die dann unter den Posts veröffentlicht wurden. Die Anweisung an die Sprachmodelle lautete, den Veröffentlicher des Posts von seiner Meinung abzubringen. Eine der zentralen Forschungsfragen befasste sich damit, ob die Personalisierung von Argumenten verschiedener Sprachmodelle auf Basis von Nutzermerkmalen deren Überzeugungskraft steigern kann. Dafür wurden verschiedene Test-Accounts angelegt, die spezifische Lebenserfahrungen repräsentieren, darunter ein Betroffener von sexueller Gewalt, ein Traumaberater und ein Kritiker der Black Lives Matter-Bewegung \autocite{redditmodcmv2025}. Im Ergebnis kamen die Forschenden zu dem Schluss, dass die KI-generierten Kommentare drei bis sechsmal überzeugender waren, als von Menschen verfasste Kommentare \autocite{unizurichextendedabstract}. Die Studie wurde nach öffentlicher Kritik und nach einem Verstoß gegen die Subreddit-Regeln nicht veröffentlicht \autocite{redditmodcmv2025}. \\
Wie der Einsatz von KI-generierten Inhalten abseits des Forschungskontexts auch im politischen Kontext zur Meinungsbeeinflussung genutzt werden kann, zeigen die Geschehnisse um die Alternative für Deutschland (AfD). Mithilfe des KI-Tools „Alternita Studio“, das Sprachmodelle von OpenAI, Google und Anthropic nutzt, werden „publizierbereite Posts mit Textfeldern vor Hintergrundbildern oder Videoclips“ erzeugt \autocite{widmer2026}. Grundlage des Tools sind Nachrichtenartikel, die anhand eines „Viralitätsscores“ gerankt werden. Auf Basis dieser Artikel und dem „persönlichen“ Profil des Anwenders werden dann automatisiert Beiträge für Soziale Medien generiert \autocite{alternitastudiohp}. Ziel ist dabei, durch die Orientierung an polarisierenden Nachrichtenartikeln bzw. den darin behandelten Themen, eine möglichst große Viralität der erzeugten Beiträge zu erreichen \autocite{widmer2026}. Die Anwälte der AfD erklärten, dass ein Export von KI-generierten Bildern aus dem Tool nicht möglich sei, ohne eine explizite Kennzeichnung der Herkunft. Allerdings stellte sich heraus, dass diese Schranke vergleichsweise einfach umgangen werden kann. KI-generierte Texte werden überhaupt nicht als solche kenntlich gemacht \autocite{correctiv2026}. Sowohl Google untersagt in seinen Richtlinien die „Umgehung von Schutzmaßnahmen gegen Missbrauch oder Sicherheitsfiltern, z. B. durch Manipulation des Modells“ \autocite{googlerichtlinie}, ebenso wie der EU AI Act \autocite{EU_AI_Act_2024, Art. 50}. \\
Warum aber eine zuverlässige Erkennung von KI-generiertem Text relevant ist, begründen \autocite{bevendorff2026overviewpan2026voightkampff}:
„Recognizing the ‚fingerprint‘ of AI text generation is the foundation for a healthy information
ecosystem in the future, but also to prevent model collapse when LLMs are trained increasingly
on their own output. However, LLM detection remains a challenge, especially across text
domains, but also in the face of potential LLM style obfuscations.” \\
Die Relevanz der Erkennung KI-generierter Texte lässt sich anhand verschiedener Dimensionen erklären. Zum einen hat sie eine Bedeutung für das Sozialgefüge. Das Weltwirtschaftsforum sieht in seinem Global Risks Report in Technologien der künstlichen Intelligenz eine Gefahr für den gesellschaftlichen Zusammenhalt. In dem von ihnen durchgeführten Global Risks Perception Survey kommen sie zu dem Schluss, dass die negativen Folgen von KI-Technologien zu den größten globalen Langzeitrisiken zählen. Mis- und Desinformation, die durch KI generiert werden, stellen bereits heute ein großes Problem dar, dass sich aufgrund der technologischen Entwicklung in der Zukunft weiter verschärfen wird. Es wird zunehmend schwerer werden, zwischen echten und künstlichen Inhalten zu unterscheiden \autocite{WEF_GlobalRisks_2026}. \\


\section{Reddit als Untersuchungsumgebung}
Die Plattform Reddit etablierte sich in den letzten Jahren zu einer der reichweitenstärksten Netzwerke im Internet. Charakteristisch für Reddit ist die Unterteilung in Subreddits mit eigenen spezifischen Regeln, Normen und Subkulturen, die sich als Mikrokosmen mit eigener Dynamik verstehen lassen \autocite{proferes2021}.
Soziale Netzwerke sind für viele ihrer Nutzer eine Art „Safe Space“. Reddit basiert also qua Definition auf genuin menschlichem Austausch, den Nutzer aktiv suchen und fordern. Somit scheint Reddit als das „letzte Refugium für echte menschliche Erfahrungen im Netz“ \autocite{schiffer2026}.  Dennoch wird geschätzt, dass sich die Zahl der vermutlich von KI generierten Posts auf Reddit zwischen 2024 und 2025 um ca. 13 \% erhöht hat \autocite{lambert2025}.
Reddit selbst definiert keine expliziten Regeln bzw. Verbote für die Nutzung generativer künstlicher Intelligenz. Allerdings sind die Subreddits in der Lage ihre eigenen Regeln im Rahmen der Reddit-Nutzungsbedingungen festzulegen. Dazu gehört auch die Möglichkeit eines Gebots oder eines Verbots der Nutzung von generativer KI zur Erstellung von Posts oder Kommentaren, die in ihrer Restriktivität unterschiedlich sein können \autocite{redditrichtlinie}.


Reddit selbst definiert keine Regeln zu KI nutzung (Quelle nutzungsbedingungen)
Subbredits mit eigenen Regeln zu KI-Nutzung etabliert --> Studie eher gegen KI nutzung Lloyd
Genau diese Kombi, also eher gegen KI aber immer mehr Inhalte KI hat Thema Reddit für forschung so relevant gemacht --> quelle proferes

Umso wichtiger ist es also, KI-generierte Inhalte insbesondere in Form von Text auf Reddit zu erkennen. \\
Soziale Netzwerke sind für viele ihrer Nutzer eine Art „Safe Space“. Sie suchen dort gezielt nach menschlichem Austausch. Reddit scheint als das „letzte Refugium für echte menschliche Erfahrungen im Netz“ \autocite{schiffer2026}. \\





\section{Zielsetzung}

Das zentrale Ziel dieser Bachelorarbeit ist die Evaluierung verschiedener Machine Learning-Modelle zur Erkennung von KI-generiertem Text. Die dazu verwendeten Daten stammen von der Plattform Reddit aus unterschiedlichen Subreddits. Um einen dazugehörigen KI-generierten Datensatz zu erzeugen, wurde das LLM Llama-3.1-8B-Instruct verwendet. Die Anweisung lautete dabei, einen Post zum jeweiligen Subreddit, auf Basis des Titels und des Inhalts des originalen Posts zu generieren. Auf diesen beiden Datensätzen soll dann untersucht werden, inwieweit bereits trainierte Modelle in der Lage sind,  zwischen von Menschen und von KI generierten Texten zu unterscheiden. \\

\subsection{Unterabschnitt}

Dies ist ein Beispiel für einen Unterabschnitt.

\subsubsection{Unterunterabschnitt}

Dies ist ein Beispiel für einen Unterunterabschnitt.

\section{Forschungsfragen}

Dies ist ein Beispiel, wie man Forschungsfragen formatiert.

\hypertarget{rq:1}{
    \subsubsection{F\textsubscript{1}: Ist dies eine Forschungsfrage?}
}

Man kann auf Forschungsfragen wie folgt verweisen: \hyperlink{rq:1}{F\textsubscript{1}} ist ein Beispiel für eine Forschungsfrage.

\section{Kapitelverweise}

Man kann auf andere Kapitel wie folgt verweisen: Kapitel~\ref{chap:conclusion} ist das Abschlusskapitel.
